\subsection{Итоговый проект: патрулирование}

Объединим в одной программе миссию по точкам, события автопилота, камеру, сервопривод, светодиоды и аварийное завершение. Программа оформлена как класс.

\begin{codefile}[python]{python-onboard/pioneer_sdk2/21_patrol_project.py}
\end{codefile}

\paragraph{Что происходит.} Класс \verb|Patrol| хранит всё состояние (дрон, камера, сервопривод, флаги) и делит логику на методы: события, действия (\verb|fly_to|, \verb|snapshot|), сценарий (\verb|run|) и завершение (\verb|shutdown|). Два флага \verb|threading.Event|: \verb|point| --- точка достигнута, \verb|abort| --- опасное событие; обработчик опасного события выставляет оба, поэтому основной поток просыпается и сразу видит причину. Реакция на разные события разная: при \verb|LOW_VOLTAGE1| заряда хватает на возврат --- вызывается \verb|rtl()|; при \verb|LOW_VOLTAGE2| автопилот уже сам начал экстренную посадку; при \verb|SHOCK| миссия прерывается. Фото подписываются номером точки и временем и складываются в отдельную папку для каждого запуска (\texttt{patrol\_ГГГГММДД\_ЧЧММСС}).

Обратите внимание на \verb|shutdown()|: он вызывается из \verb|finally| при любом исходе и сам решает, что делать, по состоянию \verb|get_fly_state()|. Это и есть шаблон безопасной программы из начала главы, развёрнутый до целого проекта. Пауза \verb|time.sleep(8)| после \verb|rtl()| --- грубое ожидание возврата; в своих программах заменяйте её циклом по \verb|get_fly_state()| (раздел «Курс и возврат домой»).

\begin{infobox}[Что запускается в симуляторе]
Проект использует \verb|Camera|, \verb|ServoCamera|, события и \verb|cv2.imwrite|; в симуляторе Мини~2 всё это есть, но запись файлов в нём зависит от версии симулятора. Для полноценной проверки запускайте проект на дроне.
\end{infobox}

\begin{taskbox}[Развитие проекта]
Добавьте журнал полёта (пример 10) и видеозапись (пример 16). Сравнивайте снимок каждой точки с предыдущим кругом (\verb|cv2.absdiff|) и сообщайте об изменениях --- «охранная» функция. Подключите нейросеть из главы~\ref{sec:rknn}: ищите людей на снимках и при обнаружении зависайте и мигайте.
\expected{два круга по периметру, по снимку в каждой вершине в папке запуска, посадка в начальной точке; при низком заряде --- красный свет, возврат и посадка; новые функции работают, не нарушая завершения через \texttt{finally}.}
\end{taskbox}
